\begin{register}{H}{TWAIFD\_DEVICE\_ID\_VERSION\_REG}{0x0000}\label{regdesc:TWAIFDDEVICEIDVERSIONREG}
\regfield{TWAIFD\_VER\_MAJOR}{8}{24}{{0x2}}%
\regfield{TWAIFD\_VER\_MINOR}{8}{16}{{0x4}}%
\regfield{TWAIFD\_DEVICE\_ID}{16}{0}{{0xcafd}}%
\reglabel{Reset}\regnewline%
\begin{regdesc}\begin{reglist}
\label{fielddesc:TWAIFDDEVICEID}\item [TWAIFD\_DEVICE\_ID] 表示 CAN IP 功能是否正确映射到其基地址。(RO)
\label{fielddesc:TWAIFDVERMINOR}\item [TWAIFD\_VER\_MINOR] 表示 TWAI FD IP 小版本号。(RO)
\label{fielddesc:TWAIFDVERMAJOR}\item [TWAIFD\_VER\_MAJOR] 表示 TWAI FD IP 大版本号。(RO)
\end{reglist}\end{regdesc}
\end{register}


\begin{register}{H}{TWAIFD\_MODE\_SETTINGS\_REG}{0x0004}\label{regdesc:TWAIFDMODESETTINGSREG}
\regfield{(reserved)}{4}{28}{0000}%
\regfield{TWAIFD\_PCHKE}{1}{27}{{0}}%
\regfield{TWAIFD\_FDRF}{1}{26}{{0}}%
\regfield{TWAIFD\_TBFBO}{1}{25}{{1}}%
\regfield{TWAIFD\_PEX}{1}{24}{{0}}%
\regfield{TWAIFD\_NISOFD}{1}{23}{{0}}%
\regfield{TWAIFD\_ENA}{1}{22}{{0}}%
\regfield{TWAIFD\_ILBP}{1}{21}{{0}}%
\regfield{TWAIFD\_RTRTH}{4}{17}{{0x0}}%
\regfield{TWAIFD\_RTRLE}{1}{16}{{0}}%
\regfield{(reserved)}{4}{12}{0000}%
\regfield{TWAIFD\_SAM}{1}{11}{{0}}%
\regfield{TWAIFD\_TXBBM}{1}{10}{{0}}%
\regfield{TWAIFD\_RXBAM}{1}{9}{{1}}%
\regfield{TWAIFD\_TSTM}{1}{8}{{0}}%
\regfield{TWAIFD\_ACF}{1}{7}{{0}}%
\regfield{TWAIFD\_ROM}{1}{6}{{0}}%
\regfield{TWAIFD\_TTTM}{1}{5}{{0}}%
\regfield{TWAIFD\_FDE}{1}{4}{{1}}%
\regfield{TWAIFD\_AFM}{1}{3}{{0}}%
\regfield{TWAIFD\_STM}{1}{2}{{0}}%
\regfield{TWAIFD\_BMM}{1}{1}{{0}}%
\regfield{TWAIFD\_RST}{1}{0}{{0}}%
\reglabel{Reset}\regnewline%
\begin{regdesc}\begin{reglist}
\label{fielddesc:TWAIFDRST}\item [TWAIFD\_RST] 配置是否软件复位。写 1 会复位 \modulename。写 1 后不需要再写 0，该字段会自动清除。\\
0：无效\\
1：复位\\(WO)
\label{fielddesc:TWAIFDBMM}\item [TWAIFD\_BMM] 配置是否启用总线监控模式。\\
0：禁用\\
1：启用\\(R/W)
\label{fielddesc:TWAIFDSTM}\item [TWAIFD\_STM] 配置是否启用自检模式。\\
0：禁用\\
1：启用\\(R/W)
\label{fielddesc:TWAIFDAFM}\item [TWAIFD\_AFM] 配置是否启用接受过滤器模式。\\
0：禁用\\
1：启用\\(R/W)
\label{fielddesc:TWAIFDFDE}\item [TWAIFD\_FDE] 配置是否启用灵活数据速率。启用后，\modulename 能够识别 CAN FD 帧（即 FDF 位为 1）。\\
0：禁用\\
1：启用\\(R/W)
\label{fielddesc:TWAIFDTTTM}\item [TWAIFD\_TTTM] 配置是否启用定时发送模式。\\
0：禁用\\
1：启用\\(R/W)
\label{fielddesc:TWAIFDROM}\item [TWAIFD\_ROM] 配置是否启用受限操作模式。\\
0：禁用\\
1：启用\\(R/W)
\label{fielddesc:TWAIFDACF}\item [TWAIFD\_ACF] 配置是否启用禁止确认模式。\\
0：禁用\\
1：启用\\(R/W)

\item [见下页...]
\end{reglist}\end{regdesc}
\end{register}

\addtocounter{Regfloat}{-1}
\begin{register}{H}{TWAIFD\_MODE\_SETTINGS\_REG}{0x0004}
\begin{regdesc}\begin{reglist}
\item [...接上页]

\label{fielddesc:TWAIFDTSTM}\item [TWAIFD\_TSTM] 配置是否启用测试模式。\\
0：禁用\\
1：启用\\(R/W)
\label{fielddesc:TWAIFDRXBAM}\item [TWAIFD\_RXBAM] 配置是否启用 RX buffer 自动模式。\\
0：禁用\\
1：启用\\(R/W)
\label{fielddesc:TWAIFDTXBBM}\item [TWAIFD\_TXBBM] 配置是否启用 TX buffer 备份模式。\\
0：禁用\\
1：启用\\(R/W)
\label{fielddesc:TWAIFDSAM}\item [TWAIFD\_SAM] 配置是否启用自主确认模式。\\
0：禁用\\
1：启用\\(R/W)
\label{fielddesc:TWAIFDRTRLE}\item [TWAIFD\_RTRLE] 配置是否启用重发限制。\\
0：禁用\\
1：启用\\(R/W)
\label{fielddesc:TWAIFDRTRTH}\item [TWAIFD\_RTRTH] 配置重发限制阈值。\hyperref[fielddesc:TWAIFDRTRLE]{TWAIFD\_RTRLE} 启用时，最大重发次数。(R/W)
\label{fielddesc:TWAIFDILBP}\item [TWAIFD\_ILBP] 配置是否启用回环模式。\\
0：禁用\\
1：启用\\(R/W)
\label{fielddesc:TWAIFDENA}\item [TWAIFD\_ENA] 配置是否启用 \modulename。\\
0：禁用\\
1：启用\\(R/W)
\label{fielddesc:TWAIFDNISOFD}\item [TWAIFD\_NISOFD] 配置 \modulename 是否符合 ISO 协议。仅在 \hyperref[fielddesc:TWAIFDENA]{TWAIFD\_ENA} = 0 时修改。\\
0：符合 ISO\\
1：符合 CAN FD 1.0\\(R/W)
\label{fielddesc:TWAIFDPEX}\item [TWAIFD\_PEX] 配置是否启用协议异常处理。仅在 \hyperref[fielddesc:TWAIFDENA]{TWAIFD\_ENA} = 0 时修改。\\
0：禁用\\
1：启用\\(R/W)

\item [见下页...]
\end{reglist}\end{regdesc}
\end{register}

\addtocounter{Regfloat}{-1}
\begin{register}{H}{TWAIFD\_MODE\_SETTINGS\_REG}{0x0004}
\begin{regdesc}\begin{reglist}
\item [...接上页]

\label{fielddesc:TWAIFDTBFBO}\item [TWAIFD\_TBFBO] 配置 \modulename 离线时，所有 TX buffer 是否进入 TX failed 状态。\\
0：不进入 TX failed\\
1：进入 TX failed\\(R/W)
\label{fielddesc:TWAIFDFDRF}\item [TWAIFD\_FDRF] 配置帧过滤器是否丢弃 RTR 帧。\\
0：接受\\
1：丢弃\\(R/W)
\label{fielddesc:TWAIFDPCHKE}\item [TWAIFD\_PCHKE] 配置是否启用 TX buffer 和 RX buffer 的奇偶校验检查。\\
0：禁用\\
1：启用\\(R/W)

\end{reglist}\end{regdesc}
\end{register}


\begin{register}{H}{TWAIFD\_COMMAND\_REG}{0x000C}\label{regdesc:TWAIFDCOMMANDREG}
\regfield{(reserved)}{21}{11}{000000000000000000000}%
\regfield{TWAIFD\_CTXDPE}{1}{10}{{0}}%
\regfield{TWAIFD\_CTXPE}{1}{9}{{0}}%
\regfield{TWAIFD\_CRXPE}{1}{8}{{0}}%
\regfield{TWAIFD\_CPEXS}{1}{7}{{0}}%
\regfield{TWAIFD\_TXFCRST}{1}{6}{{0}}%
\regfield{TWAIFD\_RXFCRST}{1}{5}{{0}}%
\regfield{TWAIFD\_ERCRST}{1}{4}{{0}}%
\regfield{TWAIFD\_CDO}{1}{3}{{0}}%
\regfield{TWAIFD\_RRB}{1}{2}{{0}}%
\regfield{TWAIFD\_RXRPMV}{1}{1}{{0}}%
\regfield{(reserved)}{1}{0}{0}%
\reglabel{Reset}\regnewline%
\begin{regdesc}\begin{reglist}
\label{fielddesc:TWAIFDRXRPMV}\item [TWAIFD\_RXRPMV] 配置是否将 RX buffer 读取指针向前移动。 \\
0：无操作\\
1：向前移动\\
(WO)
\label{fielddesc:TWAIFDRRB}\item [TWAIFD\_RRB] 配置是否刷新 RX buffer 并重置其内存指针。\\
0：无效\\
1：刷新\\(WO)
\label{fielddesc:TWAIFDCDO}\item [TWAIFD\_CDO] 配置是否清除 RX buffer 的数据溢出标志。\\
0：无效\\
1：清除\\(WO)
\label{fielddesc:TWAIFDERCRST}\item [TWAIFD\_ERCRST] 配置是否复位错误计数器。若 \modulename{} 不处于离线状态，或因禁用（\hyperref[fielddesc:TWAIFDENA]{TWAIFD\_ENA} = 0）导致离线，此命令无效。\\
0：无效\\
1：复位\\(WO)
\label{fielddesc:TWAIFDRXFCRST}\item [TWAIFD\_RXFCRST] 配置是否清除 RX 总线流量计数器 RX\_COUNTER。\\
0：无效\\
1：清除\\(WO)
\label{fielddesc:TWAIFDTXFCRST}\item [TWAIFD\_TXFCRST] 配置是否清除 TX 总线流量计数器 TX\_COUNTER。\\
0：无效\\
1：清除\\(WO)
\label{fielddesc:TWAIFDCPEXS}\item [TWAIFD\_CPEXS] 配置是否清除协议异常标志。\\
0：无效\\
1：清除\\(WO)

\item [见下页...]
\end{reglist}\end{regdesc}
\end{register}

\addtocounter{Regfloat}{-1}
\begin{register}{H}{TWAIFD\_COMMAND\_REG}{0x000C}
\begin{regdesc}\begin{reglist}
\item [...接上页]

\label{fielddesc:TWAIFDCRXPE}\item [TWAIFD\_CRXPE] 配置是否清除 \hyperref[fielddesc:TWAIFDRXPE]{TWAIFD\_RXPE} 标志。\\
0：无效\\
1：清除\\(WO)
\label{fielddesc:TWAIFDCTXPE}\item [TWAIFD\_CTXPE] 配置是否清除 \hyperref[fielddesc:TWAIFDTXPE]{TWAIFD\_TXPE} 标志。\\
0：无效\\
1：清除\\(WO)
\label{fielddesc:TWAIFDCTXDPE}\item [TWAIFD\_CTXDPE] 配置是否清除双重奇偶校验错误标志。\\
0：无效\\
1：清除\\(WO)
\end{reglist}\end{regdesc}
\end{register}


\begin{register}{H}{TWAIFD\_INT\_STAT\_REG}{0x0010}\label{regdesc:TWAIFDINTSTATREG}
\regfield{(reserved)}{20}{12}{00000000000000000000}%
\regfield{TWAIFD\_TXBHCI\_INT\_ST}{1}{11}{{0}}%
\regfield{TWAIFD\_RBNEI\_INT\_ST}{1}{10}{{0}}%
\regfield{TWAIFD\_BSI\_INT\_ST}{1}{9}{{0}}%
\regfield{TWAIFD\_RXFI\_INT\_ST}{1}{8}{{0}}%
\regfield{TWAIFD\_OFI\_INT\_ST}{1}{7}{{0}}%
\regfield{TWAIFD\_BEI\_INT\_ST}{1}{6}{{0}}%
\regfield{TWAIFD\_ALI\_INT\_ST}{1}{5}{{0}}%
\regfield{TWAIFD\_FCSI\_INT\_ST}{1}{4}{{0}}%
\regfield{TWAIFD\_DOI\_INT\_ST}{1}{3}{{0}}%
\regfield{TWAIFD\_EWLI\_INT\_ST}{1}{2}{{0}}%
\regfield{TWAIFD\_TXI\_INT\_ST}{1}{1}{{0}}%
\regfield{TWAIFD\_RXI\_INT\_ST}{1}{0}{{0}}%
\reglabel{Reset}\regnewline%
\begin{regdesc}\begin{reglist}
\label{fielddesc:TWAIFDRXIINTST}\item [TWAIFD\_RXI\_INT\_ST] TWAIFD\_RXI\_INT 的屏蔽中断状态。(R/W1C)
\label{fielddesc:TWAIFDTXIINTST}\item [TWAIFD\_TXI\_INT\_ST] TWAIFD\_TXI\_INT 的屏蔽中断状态。(R/W1C)
\label{fielddesc:TWAIFDEWLIINTST}\item [TWAIFD\_EWLI\_INT\_ST] TWAIFD\_EWLI\_INT 的屏蔽中断状态。(R/W1C)
\label{fielddesc:TWAIFDDOIINTST}\item [TWAIFD\_DOI\_INT\_ST] TWAIFD\_DOI\_INT 的屏蔽中断状态。(R/W1C)
\label{fielddesc:TWAIFDFCSIINTST}\item [TWAIFD\_FCSI\_INT\_ST] TWAIFD\_FCSI\_INT 的屏蔽中断状态。(R/W1C)
\label{fielddesc:TWAIFDALIINTST}\item [TWAIFD\_ALI\_INT\_ST] TWAIFD\_ALI\_INT 的屏蔽中断状态。(R/W1C)
\label{fielddesc:TWAIFDBEIINTST}\item [TWAIFD\_BEI\_INT\_ST] TWAIFD\_BEI\_INT 的屏蔽中断状态。(R/W1C)
\label{fielddesc:TWAIFDOFIINTST}\item [TWAIFD\_OFI\_INT\_ST] TWAIFD\_OFI\_INT 的屏蔽中断状态。(R/W1C)
\label{fielddesc:TWAIFDRXFIINTST}\item [TWAIFD\_RXFI\_INT\_ST] TWAIFD\_RXFI\_INT 的屏蔽中断状态。(R/W1C)
\label{fielddesc:TWAIFDBSIINTST}\item [TWAIFD\_BSI\_INT\_ST] TWAIFD\_BSI\_INT 的屏蔽中断状态。(R/W1C)
\label{fielddesc:TWAIFDRBNEIINTST}\item [TWAIFD\_RBNEI\_INT\_ST] TWAIFD\_RBNEI\_INT 的屏蔽中断状态。(R/W1C)
\label{fielddesc:TWAIFDTXBHCIINTST}\item [TWAIFD\_TXBHCI\_INT\_ST] TWAIFD\_TXBHCI\_INT 的屏蔽中断状态。(R/W1C)
\end{reglist}\end{regdesc}
\end{register}


\begin{register}{H}{TWAIFD\_BTR\_REG}{0x0024}\label{regdesc:TWAIFDBTRREG}
\regfield{TWAIFD\_SJW}{5}{27}{{0x2}}%
\regfield{TWAIFD\_BRP}{8}{19}{{0xa}}%
\regfield{TWAIFD\_PH2}{6}{13}{{0x5}}%
\regfield{TWAIFD\_PH1}{6}{7}{{0x3}}%
\regfield{TWAIFD\_PROP}{7}{0}{{0x5}}%
\reglabel{Reset}\regnewline%
\begin{regdesc}\begin{reglist}
\label{fielddesc:TWAIFDPROP}\item [TWAIFD\_PROP] 配置标称位速率的传播段。\\
测量单位：时间定额。(R/W)
\label{fielddesc:TWAIFDPH1}\item [TWAIFD\_PH1] 配置标称位速率的相位 1 段。\\
测量单位：时间定额。(R/W)
\label{fielddesc:TWAIFDPH2}\item [TWAIFD\_PH2] 配置标称位速率的相位 2 段。\\
测量单位：时间定额。(R/W)
\label{fielddesc:TWAIFDBRP}\item [TWAIFD\_BRP] 配置标称位速率的波特率预分频器。\\
测量单位：系统时钟周期。(R/W)
\label{fielddesc:TWAIFDSJW}\item [TWAIFD\_SJW] 配置标称位时间内的同步跳宽。\\
测量单位：时间定额。(R/W)
\end{reglist}\end{regdesc}
\end{register}


\begin{register}{H}{TWAIFD\_BTR\_FD\_REG}{0x0028}\label{regdesc:TWAIFDBTRFDREG}
\regfield{TWAIFD\_SJW\_FD}{5}{27}{{0x2}}%
\regfield{TWAIFD\_BRP\_FD}{8}{19}{{0x4}}%
\regfield{(reserved)}{1}{18}{0}%
\regfield{TWAIFD\_PH2\_FD}{5}{13}{{0x3}}%
\regfield{(reserved)}{1}{12}{0}%
\regfield{TWAIFD\_PH1\_FD}{5}{7}{{0x3}}%
\regfield{(reserved)}{1}{6}{0}%
\regfield{TWAIFD\_PROP\_FD}{6}{0}{{0x3}}%
\reglabel{Reset}\regnewline%
\begin{regdesc}\begin{reglist}
\label{fielddesc:TWAIFDPROPFD}\item [TWAIFD\_PROP\_FD] 配置数据位速率的传播段。\\
测量单位：时间定额。(R/W)
\label{fielddesc:TWAIFDPH1FD}\item [TWAIFD\_PH1\_FD] 配置数据位速率的相位 1 段。\\
测量单位：时间定额。(R/W)
\label{fielddesc:TWAIFDPH2FD}\item [TWAIFD\_PH2\_FD] 配置数据位速率的相位 2 段。\\
测量单位：时间定额。(R/W)
\label{fielddesc:TWAIFDBRPFD}\item [TWAIFD\_BRP\_FD] 配置数据位速率的波特率预分频器。\\
测量单位：核心时钟周期。(R/W)
\label{fielddesc:TWAIFDSJWFD}\item [TWAIFD\_SJW\_FD] 配置数据位时间内的同步跳宽。\\
测量单位：时间定额。(R/W)
\end{reglist}\end{regdesc}
\end{register}


\begin{register}{H}{TWAIFD\_TRV\_DELAY\_SSP\_CFG\_REG}{0x0080}\label{regdesc:TWAIFDTRVDELAYSSPCFGREG}
\regfield{(reserved)}{6}{26}{000000}%
\regfield{TWAIFD\_SSP\_SRC}{2}{24}{{0x0}}%
\regfield{TWAIFD\_SSP\_OFFSET}{8}{16}{{0xa}}%
\regfield{(reserved)}{9}{7}{000000000}%
\regfield{TWAIFD\_TRV\_DELAY\_VALUE}{7}{0}{{0x0}}%
\reglabel{Reset}\regnewline%
\begin{regdesc}\begin{reglist}
\label{fielddesc:TWAIFDTRVDELAYVALUE}\item [TWAIFD\_TRV\_DELAY\_VALUE] 表示以最小时间定额的倍数测量的发送器延迟。(RO)
\label{fielddesc:TWAIFDSSPOFFSET}\item [TWAIFD\_SSP\_OFFSET] 表示以最小时间定额的倍数测量的二次采样点偏移。(R/W)
\label{fielddesc:TWAIFDSSPSRC}\item [TWAIFD\_SSP\_SRC] 表示二次采样点来源。\\
0：SSP\_SRC\_MEAS\_N\_OFFSET，SSP 位置 = TRV\_DELAY（测量的发送器延迟）+ SSP\_OFFSET。\\
1：SSP\_SRC\_NO\_SSP，不使用 SSP。发送器在数据位速率阶段使用常规采样点。\\
2：SSP\_SRC\_OFFSET，SSP 位置 = SSP\_OFFSET。忽略测量的发送器延延迟。 \\(R/W)
\end{reglist}\end{regdesc}
\end{register}


\begin{register}{H}{TWAIFD\_TIMER\_CLK\_EN\_REG}{0x0FD4}\label{regdesc:TWAIFDTIMERCLKENREG}
\regfield{(reserved)}{30}{2}{000000000000000000000000000000}%
\regfield{TWAIFD\_FORCE\_RXBUF\_MEM\_CLK\_ON}{1}{1}{{0}}%
\regfield{TWAIFD\_CLK\_EN}{1}{0}{{0}}%
\reglabel{Reset}\regnewline%
\begin{regdesc}\begin{reglist}
\label{fielddesc:TWAIFDCLKEN}\item [TWAIFD\_CLK\_EN] 配置是否强制启用寄存器配置时钟信号。\\
0：禁用\\
1：启用\\(R/W)
\label{fielddesc:TWAIFDFORCERXBUFMEMCLKON}\item [TWAIFD\_FORCE\_RXBUF\_MEM\_CLK\_ON] 配置是否强制启用 RX buffer RAM 时钟信号。\\
0：禁用\\
1：启用\\(R/W)
\end{reglist}\end{regdesc}
\end{register}


\begin{register}{H}{TWAIFD\_TIMER\_INT\_RAW\_REG}{0x0FD8}\label{regdesc:TWAIFDTIMERINTRAWREG}
\regfield{(reserved)}{31}{1}{0000000000000000000000000000000}%
\regfield{TWAIFD\_TIMER\_OVERFLOW\_INT\_RAW}{1}{0}{{0}}%
\reglabel{Reset}\regnewline%
\begin{regdesc}\begin{reglist}
\label{fielddesc:TWAIFDTIMEROVERFLOWINTRAW}\item [TWAIFD\_TIMER\_OVERFLOW\_INT\_RAW] 定时器溢出中断的原始中断状态。(R/SS/WTC)
\end{reglist}\end{regdesc}
\end{register}


\begin{register}{H}{TWAIFD\_TIMER\_INT\_ST\_REG}{0x0FDC}\label{regdesc:TWAIFDTIMERINTSTREG}
\regfield{(reserved)}{31}{1}{0000000000000000000000000000000}%
\regfield{TWAIFD\_TIMER\_OVERFLOW\_INT\_ST}{1}{0}{{0}}%
\reglabel{Reset}\regnewline%
\begin{regdesc}\begin{reglist}
\label{fielddesc:TWAIFDTIMEROVERFLOWINTST}\item [TWAIFD\_TIMER\_OVERFLOW\_INT\_ST] 定时器溢出中断的屏蔽中断状态。(RO)
\end{reglist}\end{regdesc}
\end{register}


\begin{register}{H}{TWAIFD\_TIMER\_INT\_ENA\_REG}{0x0FE0}\label{regdesc:TWAIFDTIMERINTENAREG}
\regfield{(reserved)}{31}{1}{0000000000000000000000000000000}%
\regfield{TWAIFD\_TIMER\_OVERFLOW\_INT\_ENA}{1}{0}{{0}}%
\reglabel{Reset}\regnewline%
\begin{regdesc}\begin{reglist}
\label{fielddesc:TWAIFDTIMEROVERFLOWINTENA}\item [TWAIFD\_TIMER\_OVERFLOW\_INT\_ENA] 写 1 启用定时器溢出中断。(R/W)
\end{reglist}\end{regdesc}
\end{register}


\begin{register}{H}{TWAIFD\_TIMER\_INT\_CLR\_REG}{0x0FE4}\label{regdesc:TWAIFDTIMERINTCLRREG}
\regfield{(reserved)}{31}{1}{0000000000000000000000000000000}%
\regfield{TWAIFD\_TIMER\_OVERFLOW\_INT\_CLR}{1}{0}{{0}}%
\reglabel{Reset}\regnewline%
\begin{regdesc}\begin{reglist}
\label{fielddesc:TWAIFDTIMEROVERFLOWINTCLR}\item [TWAIFD\_TIMER\_OVERFLOW\_INT\_CLR] 写 1 清除定时器溢出中断。(WT)
\end{reglist}\end{regdesc}
\end{register}


\begin{register}{H}{TWAIFD\_TIMER\_CFG\_REG}{0x0FE8}\label{regdesc:TWAIFDTIMERCFGREG}
\regfield{TWAIFD\_TIMER\_STEP}{16}{16}{{0x00}}%
\regfield{(reserved)}{7}{9}{0000000}%
\regfield{TWAIFD\_TIMER\_UP\_DN}{1}{8}{{1}}%
\regfield{(reserved)}{5}{3}{00000}%
\regfield{TWAIFD\_TIMER\_SET}{1}{2}{{0}}%
\regfield{TWAIFD\_TIMER\_CLR}{1}{1}{{0}}%
\regfield{TWAIFD\_TIMER\_CE}{1}{0}{{0}}%
\reglabel{Reset}\regnewline%
\begin{regdesc}\begin{reglist}
\label{fielddesc:TWAIFDTIMERCE}\item [TWAIFD\_TIMER\_CE] 配置是否启用定时器。\\
0：禁用\\
1：启用\\(R/W)
\label{fielddesc:TWAIFDTIMERCLR}\item [TWAIFD\_TIMER\_CLR] 配置是否清除定时器。\\
0：不清除\\
1：清除\\(WT)
\label{fielddesc:TWAIFDTIMERSET}\item [TWAIFD\_TIMER\_SET] 配置是否设置定时器。\\
0：不设置\\
1：设置值为 ld\_val\\(WT)
\label{fielddesc:TWAIFDTIMERUPDN}\item [TWAIFD\_TIMER\_UP\_DN] 配置定时器向上或向下计数。\\
0：向下\\
1：向上\\(R/W)
\label{fielddesc:TWAIFDTIMERSTEP}\item [TWAIFD\_TIMER\_STEP] 配置定时器计数步长，步长 = TWAIFD\_TIMER\_STEP + 1。(R/W)
\end{reglist}\end{regdesc}
\end{register}


\begin{register}{H}{TWAIFD\_TIMER\_LD\_VAL\_L\_REG}{0x0FEC}\label{regdesc:TWAIFDTIMERLDVALLREG}
\regfield{TWAIFD\_TIMER\_LD\_VAL\_L}{32}{0}{{0x000000}}%
\reglabel{Reset}\regnewline%
\begin{regdesc}\begin{reglist}
\label{fielddesc:TWAIFDTIMERLDVALL}\item [TWAIFD\_TIMER\_LD\_VAL\_L] 配置定时器初始值的低位。(R/W)
\end{reglist}\end{regdesc}
\end{register}


\begin{register}{H}{TWAIFD\_TIMER\_LD\_VAL\_H\_REG}{0x0FF0}\label{regdesc:TWAIFDTIMERLDVALHREG}
\regfield{TWAIFD\_TIMER\_LD\_VAL\_H}{32}{0}{{0x000000}}%
\reglabel{Reset}\regnewline%
\begin{regdesc}\begin{reglist}
\label{fielddesc:TWAIFDTIMERLDVALH}\item [TWAIFD\_TIMER\_LD\_VAL\_H] 配置定时器初始值的高位。如果时间戳有效位宽小于 33，则此字段被忽略。(R/W)
\end{reglist}\end{regdesc}
\end{register}


\begin{register}{H}{TWAIFD\_TIMER\_CT\_VAL\_L\_REG}{0x0FF4}\label{regdesc:TWAIFDTIMERCTVALLREG}
\regfield{TWAIFD\_TIMER\_CT\_VAL\_L}{32}{0}{{0xffffffff}}%
\reglabel{Reset}\regnewline%
\begin{regdesc}\begin{reglist}
\label{fielddesc:TWAIFDTIMERCTVALL}\item [TWAIFD\_TIMER\_CT\_VAL\_L] 配置定时器终到值的低位。(R/W)
\end{reglist}\end{regdesc}
\end{register}


\begin{register}{H}{TWAIFD\_TIMER\_CT\_VAL\_H\_REG}{0x0FF8}\label{regdesc:TWAIFDTIMERCTVALHREG}
\regfield{TWAIFD\_TIMER\_CT\_VAL\_H}{32}{0}{{0xffffffff}}%
\reglabel{Reset}\regnewline%
\begin{regdesc}\begin{reglist}
\label{fielddesc:TWAIFDTIMERCTVALH}\item [TWAIFD\_TIMER\_CT\_VAL\_H] 配置定时器终到值的高位。如果时间戳有效位宽小于 33，则此字段被忽略。(R/W)
\end{reglist}\end{regdesc}
\end{register}


\begin{register}{H}{TWAIFD\_STATUS\_REG}{0x0008}\label{regdesc:TWAIFDSTATUSREG}
\regfield{(reserved)}{14}{18}{0000000000000}%
\regfield{TWAIFD\_STRGS}{1}{17}{{1}}%
\regfield{TWAIFD\_STCNT}{1}{16}{{1}}%
\regfield{(reserved)}{4}{12}{0000}%
\regfield{TWAIFD\_TXDPE}{1}{11}{{0}}%
\regfield{TWAIFD\_TXPE}{1}{10}{{0}}%
\regfield{TWAIFD\_RXPE}{1}{9}{{0}}%
\regfield{TWAIFD\_PEXS}{1}{8}{{0}}%
\regfield{TWAIFD\_IDLE}{1}{7}{{1}}%
\regfield{TWAIFD\_EWL}{1}{6}{{0}}%
\regfield{TWAIFD\_TXS}{1}{5}{{0}}%
\regfield{TWAIFD\_RXS}{1}{4}{{0}}%
\regfield{TWAIFD\_EFT}{1}{3}{{0}}%
\regfield{TWAIFD\_TXNF}{1}{2}{{1}}%
\regfield{TWAIFD\_DOR}{1}{1}{{0}}%
\regfield{TWAIFD\_RXNE}{1}{0}{{0}}%
\reglabel{Reset}\regnewline%
\begin{regdesc}\begin{reglist}
\label{fielddesc:TWAIFDRXNE}\item [TWAIFD\_RXNE] 表示 RX buffer 不为空。当 buffer 中存储至少一帧数据时，该字段为 1。\\
0：为空\\
1：不为空\\(RO)
\label{fielddesc:TWAIFDDOR}\item [TWAIFD\_DOR] 表示数据溢出标志。当由于 RX buffer 空间不足导致帧被丢弃时，该字段被置位。可通过 \hyperref[fielddesc:TWAIFDRRB]{TWAIFD\_RRB} 清除。\\
0：未溢出\\
1：溢出\\(RO)
\label{fielddesc:TWAIFDTXNF}\item [TWAIFD\_TXNF] 表示 TX buffer 的状态。当至少有一个 TX buffer 为空时，该字段被置位。\\
0：未满\\
1：已满\\(RO)
\label{fielddesc:TWAIFDEFT}\item [TWAIFD\_EFT] 表示是否正在发送错误帧。\\
0：未发送\\
1：正在发送\\(RO)
\label{fielddesc:TWAIFDRXS}\item [TWAIFD\_RXS] 表示 \modulename 是否为 CAN 帧的接收方。\\
0：未接收\\
1：正在接收\\(RO)
\label{fielddesc:TWAIFDTXS}\item [TWAIFD\_TXS] 表示 \modulename 是否为 CAN 帧的发送方。\\
0：未发送\\
1：正在发送\\(RO)
\label{fielddesc:TWAIFDEWL}\item [TWAIFD\_EWL] 表示 TEC 或 REC 是否达到（大于或等于）EWL。\\
0：未达到\\
1：达到\\(RO)
\label{fielddesc:TWAIFDIDLE}\item [TWAIFD\_IDLE] 表示是否总线处于空闲状态（不在发送或接收帧）或 \modulename 处于离线状态。\\
0：非空闲\\
1：空闲\\(RO)

\item [见下页...]
\end{reglist}\end{regdesc}
\end{register}

\addtocounter{Regfloat}{-1}
\begin{register}{H}{TWAIFD\_STATUS\_REG}{0x0008}
\begin{regdesc}\begin{reglist}
\item [...接上页]

\label{fielddesc:TWAIFDPEXS}\item [TWAIFD\_PEXS] 表示是否发生协议异常。\\
0：未发生\\
1：发生，可通过向 \hyperref[fielddesc:TWAIFDCPEXS]{TWAIFD\_CPEXS} 写 1 清除\\(RO)
\label{fielddesc:TWAIFDRXPE}\item [TWAIFD\_RXPE] 表示通过 \hyperref[regdesc:TWAIFDRXDATAREG]{TWAIFD\_RX\_DATA} 从 RX buffer 读取 CAN 帧时是否检测到奇偶校验错误。\\
0：未检测到\\
1：已检测到\\(RO)
\label{fielddesc:TWAIFDTXPE}\item [TWAIFD\_TXPE] 表示从 TX buffer 发送帧时是否检测到奇偶校验错误。\\
0：未检测到\\
1：已检测到\\(RO)
\label{fielddesc:TWAIFDTXDPE}\item [TWAIFD\_TXDPE] 表示在 TX buffer 备份模式下，备份 TX buffer 是否检测到双重奇偶校验错误。\\
0：未检测到\\
1：已检测到\\(RO)
\label{fielddesc:TWAIFDSTCNT}\item [TWAIFD\_STCNT] 表示是否启用流量计数器。\\
0：禁用\\
1：启用\\(RO)
\label{fielddesc:TWAIFDSTRGS}\item [TWAIFD\_STRGS] 表示是否启用用于存储器可测试性的测试寄存器。\\
0：禁用\\
1：启用\\(RO)
\end{reglist}\end{regdesc}
\end{register}

\begin{register}{H}{TWAIFD\_RX\_MEM\_INFO\_REG}{0x0060}\label{regdesc:TWAIFDRXMEMINFOREG}
\regfield{(reserved)}{3}{29}{000}%
\regfield{TWAIFD\_RX\_FREE}{13}{16}{{0x80}}%
\regfield{(reserved)}{3}{13}{000}%
\regfield{TWAIFD\_RX\_BUFF\_SIZE}{13}{0}{{0x80}}%
\reglabel{Reset}\regnewline%
\begin{regdesc}\begin{reglist}
\label{fielddesc:TWAIFDRXBUFFSIZE}\item [TWAIFD\_RX\_BUFF\_SIZE] 表示 RX buffer 的大小，以 32 位字为单位。(RO)
\label{fielddesc:TWAIFDRXFREE}\item [TWAIFD\_RX\_FREE] 表示 RX buffer 空闲区域的大小，以 32 位字为单位。(RO)
\end{reglist}\end{regdesc}
\end{register}


\begin{register}{H}{TWAIFD\_RX\_POINTERS\_REG}{0x0064}\label{regdesc:TWAIFDRXPOINTERSREG}
\regfield{(reserved)}{4}{28}{0000}%
\regfield{TWAIFD\_RX\_RPP}{12}{16}{{0x0}}%
\regfield{(reserved)}{4}{12}{0000}%
\regfield{TWAIFD\_RX\_WPP}{12}{0}{{0x0}}%
\reglabel{Reset}\regnewline%
\begin{regdesc}\begin{reglist}
\label{fielddesc:TWAIFDRXWPP}\item [TWAIFD\_RX\_WPP] 表示 RX buffer 中的写指针位置。在存储接收到的帧时，写指针会更新。(RO)
\label{fielddesc:TWAIFDRXRPP}\item [TWAIFD\_RX\_RPP] 表示 RX buffer 中的读指针位置。在读取接收到的帧时，读指针会更新。(RO)
\end{reglist}\end{regdesc}
\end{register}


\begin{register}{H}{TWAIFD\_RX\_STATUS\_RX\_SETTINGS\_REG}{0x0068}\label{regdesc:TWAIFDRXSTATUSRXSETTINGSREG}
\regfield{(reserved)}{15}{17}{000000000000000}%
\regfield{TWAIFD\_RTSOP}{1}{16}{{0}}%
\regfield{(reserved)}{1}{15}{0}%
\regfield{TWAIFD\_RXFRC}{11}{4}{{0x0}}%
\regfield{(reserved)}{1}{3}{0}%
\regfield{TWAIFD\_RXMOF}{1}{2}{{0}}%
\regfield{TWAIFD\_RXF}{1}{1}{{0}}%
\regfield{TWAIFD\_RXE}{1}{0}{{1}}%
\reglabel{Reset}\regnewline%
\begin{regdesc}\begin{reglist}
\label{fielddesc:TWAIFDRXE}\item [TWAIFD\_RXE] 表示 RX buffer 是否为空。\\
0：不为空\\
1：为空\\(RO)
\label{fielddesc:TWAIFDRXF}\item [TWAIFD\_RXF] 表示 RX buffer 是否已满。\\
0：未满\\
1：已满\\(RO)
\label{fielddesc:TWAIFDRXMOF}\item [TWAIFD\_RXMOF] 表示 RX buffer 中接收到的帧数。RXMOF 表示 RX buffer 正在读取一帧数据。当该字段为 1 时，下次从 \hyperref[regdesc:TWAIFDRXDATAREG]{TWAIFD\_RX\_DATA} 读取的数据将不是 CAN 帧的第一个字（FRAME\_FORMAT\_W）。(RO)
\label{fielddesc:TWAIFDRXFRC}\item [TWAIFD\_RXFRC] 表示 RX buffer 当前存储的 CAN 帧数量。(RO)
\label{fielddesc:TWAIFDRTSOP}\item [TWAIFD\_RTSOP] 表示 RX buffer 时间戳选项。仅在 \hyperref[fielddesc:TWAIFDENA]{TWAIFD\_ENA} = 0 时修改。\\
0：RTS\_END，RX FIFO 接收到的帧的时间戳记录在 EOF 字段的最后一位\\
1：RTS\_BEG，RX FIFO 接收到的帧的时间戳记录在 SOF 字段\\(R/W)
\end{reglist}\end{regdesc}
\end{register}


\begin{register}{H}{TWAIFD\_TX\_STATUS\_REG}{0x0070}\label{regdesc:TWAIFDTXSTATUSREG}
\regfield{TWAIFD\_TX8S}{4}{28}{{0x0}}%
\regfield{TWAIFD\_TX7S}{4}{24}{{0x0}}%
\regfield{TWAIFD\_TX6S}{4}{20}{{0x0}}%
\regfield{TWAIFD\_TX5S}{4}{16}{{0x0}}%
\regfield{TWAIFD\_TX4S}{4}{12}{{0x8}}%
\regfield{TWAIFD\_TX3S}{4}{8}{{0x8}}%
\regfield{TWAIFD\_TX2S}{4}{4}{{0x8}}%
\regfield{TWAIFD\_TXTB0\_STATE}{4}{0}{{0x8}}%
\reglabel{Reset}\regnewline%
\begin{regdesc}\begin{reglist}
\label{fielddesc:TWAIFDTXTB0STATE}\item [TWAIFD\_TXTB0\_STATE] 表示 TX buffer 1 的状态。\\
0：TXT\_NOT\_EXIST，TX buffer 在核心中不存在（适用于 \modulename 有少于 8 个 TX buffer）\\
1：TXT\_RDY，TX buffer 处于 ready 状态，等待 \modulename 启动发送\\
2：TXT\_TRAN，TX buffer 处于 TX in progress 状态，\modulename 正在发送帧\\
3：TXT\_ABTP，TX buffer 处于 abort in progress 状态\\
4：TXT\_TOK，TX buffer 处于 TX OK 状态\\
6：TXT\_ERR，TX buffer 处于 failed 状态\\
7：TXT\_ABT，TX buffer 处于 aborted 状态\\
8：TXT\_ETY，TX buffer 处于 empty 状态\\(RO)
\label{fielddesc:TWAIFDTX2S}\item [TWAIFD\_TX2S] 表示 TX buffer 2 的状态。字段含义参考 \hyperref[fielddesc:TWAIFDTXTB0STATE]{TWAIFD\_TXTB0\_STATE}。(RO)
\label{fielddesc:TWAIFDTX3S}\item [TWAIFD\_TX3S] 表示 TX buffer 3 的状态。字段含义参考 \hyperref[fielddesc:TWAIFDTXTB0STATE]{TWAIFD\_TXTB0\_STATE}。(RO)
\label{fielddesc:TWAIFDTX4S}\item [TWAIFD\_TX4S] 表示 TX buffer 4 的状态。字段含义参考 \hyperref[fielddesc:TWAIFDTXTB0STATE]{TWAIFD\_TXTB0\_STATE}。(RO)
\label{fielddesc:TWAIFDTX5S}\item [TWAIFD\_TX5S] 表示 TX buffer 5 的状态。字段含义参考 \hyperref[fielddesc:TWAIFDTXTB0STATE]{TWAIFD\_TXTB0\_STATE}。(RO)
\label{fielddesc:TWAIFDTX6S}\item [TWAIFD\_TX6S] 表示 TX buffer 6 的状态。字段含义参考 \hyperref[fielddesc:TWAIFDTXTB0STATE]{TWAIFD\_TXTB0\_STATE}。(RO)
\label{fielddesc:TWAIFDTX7S}\item [TWAIFD\_TX7S] 表示 TX buffer 7 的状态。字段含义参考 \hyperref[fielddesc:TWAIFDTXTB0STATE]{TWAIFD\_TXTB0\_STATE}。(RO)
\label{fielddesc:TWAIFDTX8S}\item [TWAIFD\_TX8S] 表示 TX buffer 8 的状态。字段含义参考 \hyperref[fielddesc:TWAIFDTXTB0STATE]{TWAIFD\_TXTB0\_STATE}。(RO)
\end{reglist}\end{regdesc}
\end{register}


\begin{register}{H}{TWAIFD\_ERR\_CAPT\_RETR\_CTR\_ALC\_TS\_INFO\_REG}{0x007C}\label{regdesc:TWAIFDERRCAPTRETRCTRALCTSINFOREG}
\regfield{(reserved)}{2}{30}{00}%
\regfield{TWAIFD\_TS\_BITS}{6}{24}{{0x0}}%
\regfield{TWAIFD\_ALC\_ID\_FIELD}{3}{21}{{0x0}}%
\regfield{TWAIFD\_ALC\_BIT}{5}{16}{{0x0}}%
\regfield{(reserved)}{4}{12}{0000}%
\regfield{TWAIFD\_RETR\_CTR\_VAL}{4}{8}{{0x0}}%
\regfield{TWAIFD\_ERR\_TYPE}{3}{5}{{0x0}}%
\regfield{TWAIFD\_ERR\_POS}{5}{0}{{0x1f}}%
\reglabel{Reset}\regnewline%
\begin{regdesc}\begin{reglist}
\label{fielddesc:TWAIFDERRPOS}\item [TWAIFD\_ERR\_POS] 表示最后一个错误的位置。 
0：ERC\_POS\_SOF，起始帧错误\\
1：ERC\_POS\_ARB，仲裁段错误\\
2：ERC\_POS\_CTRL，控制段错误\\
3：ERC\_POS\_DATA，数据段错误\\
4：ERC\_POS\_CRC，CRC 段错误\\
5：ERC\_POS\_ACK，CRC 分隔符、ACK 或 ACK 分隔符错误\\
6：ERC\_POS\_EOF，帧结束段错误\\
7：ERC\_POS\_ERR，错误帧期间出现错误\\
8：ERC\_POS\_OVRL，过载帧错误\\
31：ERC\_POS\_OTHER，其他位置错误\\(RO)
\label{fielddesc:TWAIFDERRTYPE}\item [TWAIFD\_ERR\_TYPE] 表示最后一个错误的类型。\\
0：ERC\_BIT\_ERR，位错误\\
1：ERC\_CRC\_ERR，CRC 错误\\
2：ERC\_FRM\_ERR，格式错误\\
3：ERC\_ACK\_ERR，ACK 错误\\
4：ERC\_STUF\_ERR，填充错误 \\(RO)
\label{fielddesc:TWAIFDRETRCTRVAL}\item [TWAIFD\_RETR\_CTR\_VAL] 表示重发计数器的当前值。(RO)
\label{fielddesc:TWAIFDALCBIT}\item [TWAIFD\_ALC\_BIT] 表示仲裁丢失时的位捕获位置。若该字段为 ALC\_BASE\_ID，则仲裁丢失的基本标识符位索引为 11 - ALC\_VAL。如果 ALC\_ID\_FIELD = ALC\_EXTENSION，则仲裁丢失的扩展标识符位索引为 18 - ALC\_VAL。其他值无效。(RO)
\label{fielddesc:TWAIFDALCIDFIELD}\item [TWAIFD\_ALC\_ID\_FIELD] 表示仲裁丢失时 CAN 标识符中的部分。\\
0：ALC\_RSVD，自上次复位后未丢失过仲裁\\
1：ALC\_BASE\_ID，在基本标识符期间丢失仲裁\\
2：ALC\_SRR\_RTR，在基本标识符之后的第一个位丢失仲裁（扩展帧的 SRR 位，CAN 2.0 Base 帧的 RTR 位）\\
3：ALC\_IDE，在 IDE 位期间丢失仲裁\\
4：ALC\_EXTENSION，在拓展标识符期间丢失仲裁\\
5：ALC\_RTR，在拓展标识符之后的 RTR 位期间丢失仲裁\\(RO)
\label{fielddesc:TWAIFDTSBITS}\item [TWAIFD\_TS\_BITS] 表示时间基准的有效位数减 1（0x3F = 64 位时间基准）。(RO)
\end{reglist}\end{regdesc}
\end{register}


\begin{register}{H}{TWAIFD\_TIMESTAMP\_LOW\_REG}{0x0094}\label{regdesc:TWAIFDTIMESTAMPLOWREG}
\regfield{TWAIFD\_TIMESTAMP\_LOW}{32}{0}{{0x000000}}%
\reglabel{Reset}\regnewline%
\begin{regdesc}\begin{reglist}
\label{fielddesc:TWAIFDTIMESTAMPLOW}\item [TWAIFD\_TIMESTAMP\_LOW] 表示时间基准的 0:31 位。(RO)
\end{reglist}\end{regdesc}
\end{register}


\begin{register}{H}{TWAIFD\_TIMESTAMP\_HIGH\_REG}{0x0098}\label{regdesc:TWAIFDTIMESTAMPHIGHREG}
\regfield{TWAIFD\_TIMESTAMP\_HIGH}{32}{0}{{0x000000}}%
\reglabel{Reset}\regnewline%
\begin{regdesc}\begin{reglist}
\label{fielddesc:TWAIFDTIMESTAMPHIGH}\item [TWAIFD\_TIMESTAMP\_HIGH] 表示时间基准的 32:63 位。(RO)
\end{reglist}\end{regdesc}
\end{register}


\begin{register}{H}{TWAIFD\_INT\_ENA\_SET\_REG}{0x0014}\label{regdesc:TWAIFDINTENASETREG}
\regfield{(reserved)}{20}{12}{00000000000000000000}%
\regfield{TWAIFD\_TXBHCI\_INT\_ENA\_MASK}{1}{11}{{0}}%
\regfield{TWAIFD\_RBNEI\_INT\_ENA\_MASK}{1}{10}{{0}}%
\regfield{TWAIFD\_BSI\_INT\_ENA\_MASK}{1}{9}{{0}}%
\regfield{TWAIFD\_RXFI\_INT\_ENA\_MASK}{1}{8}{{0}}%
\regfield{TWAIFD\_OFI\_INT\_ENA\_MASK}{1}{7}{{0}}%
\regfield{TWAIFD\_BEI\_INT\_ENA\_MASK}{1}{6}{{0}}%
\regfield{TWAIFD\_ALI\_INT\_ENA\_MASK}{1}{5}{{0}}%
\regfield{TWAIFD\_FCSI\_INT\_ENA\_MASK}{1}{4}{{0}}%
\regfield{TWAIFD\_DOI\_INT\_ENA\_MASK}{1}{3}{{0}}%
\regfield{TWAIFD\_EWLI\_INT\_ENA\_MASK}{1}{2}{{0}}%
\regfield{TWAIFD\_TXI\_INT\_ENA\_MASK}{1}{1}{{0}}%
\regfield{TWAIFD\_RXI\_INT\_ENA\_MASK}{1}{0}{{0}}%
\reglabel{Reset}\regnewline%
\begin{regdesc}\begin{reglist}
\label{fielddesc:TWAIFDRXIINTENAMASK}\item [TWAIFD\_RXI\_INT\_ENA\_MASK] 写 1 使能 TWAIFD\_RXI\_INT。(R/W1S)
\label{fielddesc:TWAIFDTXIINTENAMASK}\item [TWAIFD\_TXI\_INT\_ENA\_MASK] 写 1 使能 TWAIFD\_TXI\_INT。(R/W1S)
\label{fielddesc:TWAIFDEWLIINTENAMASK}\item [TWAIFD\_EWLI\_INT\_ENA\_MASK] 写 1 使能 TWAIFD\_EWLI\_INT。(R/W1S)
\label{fielddesc:TWAIFDDOIINTENAMASK}\item [TWAIFD\_DOI\_INT\_ENA\_MASK] 写 1 使能 TWAIFD\_DOI\_INT。(R/W1S)
\label{fielddesc:TWAIFDFCSIINTENAMASK}\item [TWAIFD\_FCSI\_INT\_ENA\_MASK] 写 1 使能 TWAIFD\_FCSI\_INT。(R/W1S)
\label{fielddesc:TWAIFDALIINTENAMASK}\item [TWAIFD\_ALI\_INT\_ENA\_MASK] 写 1 使能 TWAIFD\_ALI\_INT。(R/W1S)
\label{fielddesc:TWAIFDBEIINTENAMASK}\item [TWAIFD\_BEI\_INT\_ENA\_MASK] 写 1 使能 TWAIFD\_BEI\_INT。(R/W1S)
\label{fielddesc:TWAIFDOFIINTENAMASK}\item [TWAIFD\_OFI\_INT\_ENA\_MASK] 写 1 使能 TWAIFD\_OFI\_INT。(R/W1S)
\label{fielddesc:TWAIFDRXFIINTENAMASK}\item [TWAIFD\_RXFI\_INT\_ENA\_MASK] 写 1 使能 TWAIFD\_RXFI\_INT。(R/W1S)
\label{fielddesc:TWAIFDBSIINTENAMASK}\item [TWAIFD\_BSI\_INT\_ENA\_MASK] 写 1 使能 TWAIFD\_BSI\_INT。(R/W1S)
\label{fielddesc:TWAIFDRBNEIINTENAMASK}\item [TWAIFD\_RBNEI\_INT\_ENA\_MASK] 写 1 使能 TWAIFD\_RBNEI\_INT。(R/W1S)
\label{fielddesc:TWAIFDTXBHCIINTENAMASK}\item [TWAIFD\_TXBHCI\_INT\_ENA\_MASK] 写 1 使能 TWAIFD\_TXBHCI\_INT。(R/W1S)
\end{reglist}\end{regdesc}
\end{register}


\begin{register}{H}{TWAIFD\_INT\_ENA\_CLR\_REG}{0x0018}\label{regdesc:TWAIFDINTENACLRREG}
\regfield{(reserved)}{20}{12}{00000000000000000000}%
\regfield{TWAIFD\_TXBHCI\_INT\_ENA\_CLR}{1}{11}{{0}}%
\regfield{TWAIFD\_RBNEI\_INT\_ENA\_CLR}{1}{10}{{0}}%
\regfield{TWAIFD\_BSI\_INT\_ENA\_CLR}{1}{9}{{0}}%
\regfield{TWAIFD\_RXFI\_INT\_ENA\_CLR}{1}{8}{{0}}%
\regfield{TWAIFD\_OFI\_INT\_ENA\_CLR}{1}{7}{{0}}%
\regfield{TWAIFD\_BEI\_INT\_ENA\_CLR}{1}{6}{{0}}%
\regfield{TWAIFD\_ALI\_INT\_ENA\_CLR}{1}{5}{{0}}%
\regfield{TWAIFD\_FCSI\_INT\_ENA\_CLR}{1}{4}{{0}}%
\regfield{TWAIFD\_DOI\_INT\_ENA\_CLR}{1}{3}{{0}}%
\regfield{TWAIFD\_EWLI\_INT\_ENA\_CLR}{1}{2}{{0}}%
\regfield{TWAIFD\_TXI\_INT\_ENA\_CLR}{1}{1}{{0}}%
\regfield{TWAIFD\_RXI\_INT\_ENA\_CLR}{1}{0}{{0}}%
\reglabel{Reset}\regnewline%
\begin{regdesc}\begin{reglist}
\label{fielddesc:TWAIFDRXIINTENACLR}\item [TWAIFD\_RXI\_INT\_ENA\_CLR] 写 1 清除 TWAIFD\_RXI\_INT\_ENA。(WO)
\label{fielddesc:TWAIFDTXIINTENACLR}\item [TWAIFD\_TXI\_INT\_ENA\_CLR] 写 1 清除 TWAIFD\_TXI\_INT\_ENA。(WO)
\label{fielddesc:TWAIFDEWLIINTENACLR}\item [TWAIFD\_EWLI\_INT\_ENA\_CLR] 写 1 清除 TWAIFD\_EWLI\_INT\_ENA。(WO)
\label{fielddesc:TWAIFDDOIINTENACLR}\item [TWAIFD\_DOI\_INT\_ENA\_CLR] 写 1 清除 TWAIFD\_DOI\_INT\_ENA。(WO)
\label{fielddesc:TWAIFDFCSIINTENACLR}\item [TWAIFD\_FCSI\_INT\_ENA\_CLR] 写 1 清除 TWAIFD\_FCSI\_INT\_ENA。(WO)
\label{fielddesc:TWAIFDALIINTENACLR}\item [TWAIFD\_ALI\_INT\_ENA\_CLR] 写 1 清除 TWAIFD\_ALI\_INT\_ENA。(WO)
\label{fielddesc:TWAIFDBEIINTENACLR}\item [TWAIFD\_BEI\_INT\_ENA\_CLR] 写 1 清除 TWAIFD\_BEI\_INT\_ENA。(WO)
\label{fielddesc:TWAIFDOFIINTENACLR}\item [TWAIFD\_OFI\_INT\_ENA\_CLR] 写 1 清除 TWAIFD\_OFI\_INT\_ENA。(WO)
\label{fielddesc:TWAIFDRXFIINTENACLR}\item [TWAIFD\_RXFI\_INT\_ENA\_CLR] 写 1 清除 TWAIFD\_RXFI\_INT\_ENA。(WO)
\label{fielddesc:TWAIFDBSIINTENACLR}\item [TWAIFD\_BSI\_INT\_ENA\_CLR] 写 1 清除 TWAIFD\_BSI\_INT\_ENA。(WO)
\label{fielddesc:TWAIFDRBNEIINTENACLR}\item [TWAIFD\_RBNEI\_INT\_ENA\_CLR] 写 1 清除 TWAIFD\_RBNEI\_INT\_ENA。(WO)
\label{fielddesc:TWAIFDTXBHCIINTENACLR}\item [TWAIFD\_TXBHCI\_INT\_ENA\_CLR] 写 1 清除 TWAIFD\_TXBHCI\_INT\_ENA。(WO)
\end{reglist}\end{regdesc}
\end{register}


\begin{register}{H}{TWAIFD\_INT\_MASK\_SET\_REG}{0x001C}\label{regdesc:TWAIFDINTMASKSETREG}
\regfield{(reserved)}{20}{12}{00000000000000000000}%
\regfield{TWAIFD\_TXBHCI\_INT\_MASK\_SET}{1}{11}{{0}}%
\regfield{TWAIFD\_RBNEI\_INT\_MASK\_SET}{1}{10}{{0}}%
\regfield{TWAIFD\_BSI\_INT\_MASK\_SET}{1}{9}{{0}}%
\regfield{TWAIFD\_RXFI\_INT\_MASK\_SET}{1}{8}{{0}}%
\regfield{TWAIFD\_OFI\_INT\_MASK\_SET}{1}{7}{{0}}%
\regfield{TWAIFD\_BEI\_INT\_MASK\_SET}{1}{6}{{0}}%
\regfield{TWAIFD\_ALI\_INT\_MASK\_SET}{1}{5}{{0}}%
\regfield{TWAIFD\_FCSI\_INT\_MASK\_SET}{1}{4}{{0}}%
\regfield{TWAIFD\_DOI\_INT\_MASK\_SET}{1}{3}{{0}}%
\regfield{TWAIFD\_EWLI\_INT\_MASK\_SET}{1}{2}{{0}}%
\regfield{TWAIFD\_TXI\_INT\_MASK\_SET}{1}{1}{{0}}%
\regfield{TWAIFD\_RXI\_INT\_MASK\_SET}{1}{0}{{0}}%
\reglabel{Reset}\regnewline%
\begin{regdesc}\begin{reglist}
\label{fielddesc:TWAIFDRXIINTMASKSET}\item [TWAIFD\_RXI\_INT\_MASK\_SET]  写 1 屏蔽 TWAIFD\_RXI\_INT 的中断状态。(R/W1S)
\label{fielddesc:TWAIFDTXIINTMASKSET}\item [TWAIFD\_TXI\_INT\_MASK\_SET]  写 1 屏蔽 TWAIFD\_TXI\_INT 的中断状态。(R/W1S)
\label{fielddesc:TWAIFDEWLIINTMASKSET}\item [TWAIFD\_EWLI\_INT\_MASK\_SET]  写 1 屏蔽 TWAIFD\_EWLI\_INT 的中断状态。(R/W1S)
\label{fielddesc:TWAIFDDOIINTMASKSET}\item [TWAIFD\_DOI\_INT\_MASK\_SET]  写 1 屏蔽 TWAIFD\_DOI\_INT 的中断状态。(R/W1S)
\label{fielddesc:TWAIFDFCSIINTMASKSET}\item [TWAIFD\_FCSI\_INT\_MASK\_SET]  写 1 屏蔽 TWAIFD\_FCSI\_INT 的中断状态。(R/W1S)
\label{fielddesc:TWAIFDALIINTMASKSET}\item [TWAIFD\_ALI\_INT\_MASK\_SET]  写 1 屏蔽 TWAIFD\_ALI\_INT 的中断状态。(R/W1S)
\label{fielddesc:TWAIFDBEIINTMASKSET}\item [TWAIFD\_BEI\_INT\_MASK\_SET]  写 1 屏蔽 TWAIFD\_BEI\_INT 的中断状态。(R/W1S)
\label{fielddesc:TWAIFDOFIINTMASKSET}\item [TWAIFD\_OFI\_INT\_MASK\_SET]  写 1 屏蔽 TWAIFD\_OFI\_INT 的中断状态。(R/W1S)
\label{fielddesc:TWAIFDRXFIINTMASKSET}\item [TWAIFD\_RXFI\_INT\_MASK\_SET]  写 1 屏蔽 TWAIFD\_RXFI\_INT 的中断状态。(R/W1S)
\label{fielddesc:TWAIFDBSIINTMASKSET}\item [TWAIFD\_BSI\_INT\_MASK\_SET]  写 1 屏蔽 TWAIFD\_BSI\_INT 的中断状态。(R/W1S)
\label{fielddesc:TWAIFDRBNEIINTMASKSET}\item [TWAIFD\_RBNEI\_INT\_MASK\_SET]  写 1 屏蔽 TWAIFD\_RBNEI\_INT 的中断状态。(R/W1S)
\label{fielddesc:TWAIFDTXBHCIINTMASKSET}\item [TWAIFD\_TXBHCI\_INT\_MASK\_SET]  写 1 屏蔽 TWAIFD\_TXBHCI\_INT 的中断状态。(R/W1S)
\end{reglist}\end{regdesc}
\end{register}



\begin{register}{H}{TWAIFD\_INT\_MASK\_CLR\_REG}{0x0020}\label{regdesc:TWAIFDINTMASKCLRREG}
\regfield{(reserved)}{20}{12}{00000000000000000000}%
\regfield{TWAIFD\_TXBHCI\_INT\_MASK\_CLR}{1}{11}{{0}}%
\regfield{TWAIFD\_RBNEI\_INT\_MASK\_CLR}{1}{10}{{0}}%
\regfield{TWAIFD\_BSI\_INT\_MASK\_CLR}{1}{9}{{0}}%
\regfield{TWAIFD\_RXFI\_INT\_MASK\_CLR}{1}{8}{{0}}%
\regfield{TWAIFD\_OFI\_INT\_MASK\_CLR}{1}{7}{{0}}%
\regfield{TWAIFD\_BEI\_INT\_MASK\_CLR}{1}{6}{{0}}%
\regfield{TWAIFD\_ALI\_INT\_MASK\_CLR}{1}{5}{{0}}%
\regfield{TWAIFD\_FCSI\_INT\_MASK\_CLR}{1}{4}{{0}}%
\regfield{TWAIFD\_DOI\_INT\_MASK\_CLR}{1}{3}{{0}}%
\regfield{TWAIFD\_EWLI\_INT\_MASK\_CLR}{1}{2}{{0}}%
\regfield{TWAIFD\_TXI\_INT\_MASK\_CLR}{1}{1}{{0}}%
\regfield{TWAIFD\_RXI\_INT\_MASK\_CLR}{1}{0}{{0}}%
\reglabel{Reset}\regnewline%
\begin{regdesc}\begin{reglist}
\label{fielddesc:TWAIFDRXIINTMASKCLR}\item [TWAIFD\_RXI\_INT\_MASK\_CLR] 写 1 清除 TWAIFD\_RXI\_INT\_MASK\_CLR。(WO)
\label{fielddesc:TWAIFDTXIINTMASKCLR}\item [TWAIFD\_TXI\_INT\_MASK\_CLR] 写 1 清除 TWAIFD\_TXI\_INT\_MASK\_CLR。(WO)
\label{fielddesc:TWAIFDEWLIINTMASKCLR}\item [TWAIFD\_EWLI\_INT\_MASK\_CLR] 写 1 清除 TWAIFD\_EWLI\_INT\_MASK\_CLR。(WO)
\label{fielddesc:TWAIFDDOIINTMASKCLR}\item [TWAIFD\_DOI\_INT\_MASK\_CLR] 写 1 清除 TWAIFD\_DOI\_INT\_MASK\_CLR。(WO)
\label{fielddesc:TWAIFDFCSIINTMASKCLR}\item [TWAIFD\_FCSI\_INT\_MASK\_CLR] 写 1 清除 TWAIFD\_FCSI\_INT\_MASK\_CLR。(WO)
\label{fielddesc:TWAIFDALIINTMASKCLR}\item [TWAIFD\_ALI\_INT\_MASK\_CLR] 写 1 清除 TWAIFD\_ALI\_INT\_MASK\_CLR。(WO)
\label{fielddesc:TWAIFDBEIINTMASKCLR}\item [TWAIFD\_BEI\_INT\_MASK\_CLR] 写 1 清除 TWAIFD\_BEI\_INT\_MASK\_CLR。(WO)
\label{fielddesc:TWAIFDOFIINTMASKCLR}\item [TWAIFD\_OFI\_INT\_MASK\_CLR] 写 1 清除 TWAIFD\_OFI\_INT\_MASK\_CLR。(WO)
\label{fielddesc:TWAIFDRXFIINTMASKCLR}\item [TWAIFD\_RXFI\_INT\_MASK\_CLR] 写 1 清除 TWAIFD\_RXFI\_INT\_MASK\_CLR。(WO)
\label{fielddesc:TWAIFDBSIINTMASKCLR}\item [TWAIFD\_BSI\_INT\_MASK\_CLR] 写 1 清除 TWAIFD\_BSI\_INT\_MASK\_CLR。(WO)
\label{fielddesc:TWAIFDRBNEIINTMASKCLR}\item [TWAIFD\_RBNEI\_INT\_MASK\_CLR] 写 1 清除 TWAIFD\_RBNEI\_INT\_MASK\_CLR。(WO)
\label{fielddesc:TWAIFDTXBHCIINTMASKCLR}\item [TWAIFD\_TXBHCI\_INT\_MASK\_CLR] 写 1 清除 TWAIFD\_TXBHCI\_INT\_MASK\_CLR。(WO)
\end{reglist}\end{regdesc}
\end{register}


\begin{register}{H}{TWAIFD\_EWL\_ERP\_FAULT\_STATE\_REG}{0x002C}\label{regdesc:TWAIFDEWLERPFAULTSTATEREG}
\regfield{(reserved)}{13}{19}{0000000000000}%
\regfield{TWAIFD\_BOF}{1}{18}{{1}}%
\regfield{TWAIFD\_ERP}{1}{17}{{0}}%
\regfield{TWAIFD\_ERA}{1}{16}{{0}}%
\regfield{TWAIFD\_ERP\_LIMIT}{8}{8}{{128}}%
\regfield{TWAIFD\_EW\_LIMIT}{8}{0}{{96}}%
\reglabel{Reset}\regnewline%
\begin{regdesc}\begin{reglist}
\label{fielddesc:TWAIFDEWLIMIT}\item [TWAIFD\_EW\_LIMIT] 配置错误警告限制。当达到错误警告阈值时，可以产生中断。错误警告阈值表示总线干扰严重。(R/W)
\label{fielddesc:TWAIFDERPLIMIT}\item [TWAIFD\_ERP\_LIMIT] 配置被动错误限制。当任一错误计数器 (REC/TEC) 超过该值时，故障限制状态将变为被动错误状态。(R/W)
\label{fielddesc:TWAIFDERA}\item [TWAIFD\_ERA] 表示主动错误故障状态。\\
0：未出现主动错误故障状态\\
1：出现主动错误故障状态\\
(RO)
\label{fielddesc:TWAIFDERP}\item [TWAIFD\_ERP] 表示被动错误故障状态。\\
0：未出现被动错误故障状态\\
1：出现被动错误故障状态\\
(RO)
\label{fielddesc:TWAIFDBOF}\item [TWAIFD\_BOF] 表示离线故障状态。\\
0：未出现离线故障状态\\
1：出现离线故障状态\\
(RO)
\end{reglist}\end{regdesc}
\end{register}


\begin{register}{H}{TWAIFD\_REC\_TEC\_REG}{0x0030}\label{regdesc:TWAIFDRECTECREG}
\regfield{(reserved)}{7}{25}{0000000}%
\regfield{TWAIFD\_TEC\_VAL}{9}{16}{{0x0}}%
\regfield{(reserved)}{7}{9}{0000000}%
\regfield{TWAIFD\_REC\_VAL}{9}{0}{{0x0}}%
\reglabel{Reset}\regnewline%
\begin{regdesc}\begin{reglist}
\label{fielddesc:TWAIFDRECVAL}\item [TWAIFD\_REC\_VAL] 表示接收器错误计数器的值。(RO)
\label{fielddesc:TWAIFDTECVAL}\item [TWAIFD\_TEC\_VAL] 表示发送器错误计数器的值。(RO)
\end{reglist}\end{regdesc}
\end{register}


\begin{register}{H}{TWAIFD\_ERR\_NORM\_ERR\_FD\_REG}{0x0034}\label{regdesc:TWAIFDERRNORMERRFDREG}
\regfield{TWAIFD\_ERR\_FD\_VAL}{16}{16}{{0x00}}%
\regfield{TWAIFD\_ERR\_NORM\_VAL}{16}{0}{{0x00}}%
\reglabel{Reset}\regnewline%
\begin{regdesc}\begin{reglist}
\label{fielddesc:TWAIFDERRNORMVAL}\item [TWAIFD\_ERR\_NORM\_VAL] 表示标称位时间内的错误数量。(RO)
\label{fielddesc:TWAIFDERRFDVAL}\item [TWAIFD\_ERR\_FD\_VAL] 表示数据位时间内的错误数量。(RO)
\end{reglist}\end{regdesc}
\end{register}


\begin{register}{H}{TWAIFD\_CTR\_PRES\_REG}{0x0038}\label{regdesc:TWAIFDCTRPRESREG}
\regfield{(reserved)}{19}{13}{0000000000000000000}%
\regfield{TWAIFD\_EFD}{1}{12}{{0}}%
\regfield{TWAIFD\_ENORM}{1}{11}{{0}}%
\regfield{TWAIFD\_PRX}{1}{10}{{0}}%
\regfield{TWAIFD\_PTX}{1}{9}{{0}}%
\regfield{TWAIFD\_CTPV}{9}{0}{{0x0}}%
\reglabel{Reset}\regnewline%
\begin{regdesc}\begin{reglist}
\label{fielddesc:TWAIFDCTPV}\item [TWAIFD\_CTPV] 配置预设值用于设置错误计数器。(WO)
\label{fielddesc:TWAIFDPTX}\item [TWAIFD\_PTX] 配置是否将接收器错误计数器设置为预定义值。\\
0：无效\\
1：设置\\(WT)
\label{fielddesc:TWAIFDPRX}\item [TWAIFD\_PRX] 配置是否将发送器错误计数器设置为预定义值。\\
0：无效\\
1：设置\\(WT)
\label{fielddesc:TWAIFDENORM}\item [TWAIFD\_ENORM] 配置是否清除标称位时间的错误计数器。\\
0：无效\\
1：清除\\(WO)
\label{fielddesc:TWAIFDEFD}\item [TWAIFD\_EFD] 配置是否清除数据位时间的错误计数器。\\
0：无效\\
1：清除\\(WO)
\end{reglist}\end{regdesc}
\end{register}


\begin{register}{H}{TWAIFD\_FILTER\_A\_MASK\_REG}{0x003C}\label{regdesc:TWAIFDFILTERAMASKREG}
\regfield{(reserved)}{3}{29}{000}%
\regfield{TWAIFD\_BIT\_MASK\_A\_VAL}{29}{0}{{0x000000}}%
\reglabel{Reset}\regnewline%
\begin{regdesc}\begin{reglist}
\label{fielddesc:TWAIFDBITMASKAVAL}\item [TWAIFD\_BIT\_MASK\_A\_VAL] 配置过滤器 A 的屏蔽值。标识符格式与 TX 或 RX buffer 的 IDENTIFIER\_W 相同。如果过滤器 A 不存在，写入此寄存器无效，读取将返回全零。(R/W)
\end{reglist}\end{regdesc}
\end{register}


\begin{register}{H}{TWAIFD\_FILTER\_A\_VAL\_REG}{0x0040}\label{regdesc:TWAIFDFILTERAVALREG}
\regfield{(reserved)}{3}{29}{000}%
\regfield{TWAIFD\_BIT\_VAL\_A\_VAL}{29}{0}{{0x000000}}%
\reglabel{Reset}\regnewline%
\begin{regdesc}\begin{reglist}
\label{fielddesc:TWAIFDBITVALAVAL}\item [TWAIFD\_BIT\_VAL\_A\_VAL] 配置过滤器 A 的位值。标识符格式与 TX 或 RX buffer 的 IDENTIFIER\_W 相同。如果过滤器 A 不存在，写入此寄存器无效，读取将返回全零。(R/W)
\end{reglist}\end{regdesc}
\end{register}


\begin{register}{H}{TWAIFD\_FILTER\_B\_MASK\_REG}{0x0044}\label{regdesc:TWAIFDFILTERBMASKREG}
\regfield{(reserved)}{3}{29}{000}%
\regfield{TWAIFD\_BIT\_MASK\_B\_VAL}{29}{0}{{0x000000}}%
\reglabel{Reset}\regnewline%
\begin{regdesc}\begin{reglist}
\label{fielddesc:TWAIFDBITMASKBVAL}\item [TWAIFD\_BIT\_MASK\_B\_VAL] 配置过滤器 B 的屏蔽值。标识符格式与 TX 或 RX buffer 的 IDENTIFIER\_W 相同。如果过滤器 B 不存在，写入此寄存器无效，读取将返回全零。(R/W)
\end{reglist}\end{regdesc}
\end{register}


\begin{register}{H}{TWAIFD\_FILTER\_B\_VAL\_REG}{0x0048}\label{regdesc:TWAIFDFILTERBVALREG}
\regfield{(reserved)}{3}{29}{000}%
\regfield{TWAIFD\_BIT\_VAL\_B\_VAL}{29}{0}{{0x000000}}%
\reglabel{Reset}\regnewline%
\begin{regdesc}\begin{reglist}
\label{fielddesc:TWAIFDBITVALBVAL}\item [TWAIFD\_BIT\_VAL\_B\_VAL] 配置过滤器 B 的位值。标识符格式与 TX 或 RX buffer 的 IDENTIFIER\_W 相同。如果过滤器 B 不存在，写入此寄存器无效，读取将返回全零。(R/W)
\end{reglist}\end{regdesc}
\end{register}


\begin{register}{H}{TWAIFD\_FILTER\_C\_MASK\_REG}{0x004C}\label{regdesc:TWAIFDFILTERCMASKREG}
\regfield{(reserved)}{3}{29}{000}%
\regfield{TWAIFD\_BIT\_MASK\_C\_VAL}{29}{0}{{0x000000}}%
\reglabel{Reset}\regnewline%
\begin{regdesc}\begin{reglist}
\label{fielddesc:TWAIFDBITMASKCVAL}\item [TWAIFD\_BIT\_MASK\_C\_VAL] 配置过滤器 C 的屏蔽值。标识符格式与 TX 或 RX buffer 的 IDENTIFIER\_W 相同。如果过滤器 C 不存在，写入此寄存器无效，读取将返回全零。(R/W)
\end{reglist}\end{regdesc}
\end{register}


\begin{register}{H}{TWAIFD\_FILTER\_C\_VAL\_REG}{0x0050}\label{regdesc:TWAIFDFILTERCVALREG}
\regfield{(reserved)}{3}{29}{000}%
\regfield{TWAIFD\_BIT\_VAL\_C\_VAL}{29}{0}{{0x000000}}%
\reglabel{Reset}\regnewline%
\begin{regdesc}\begin{reglist}
\label{fielddesc:TWAIFDBITVALCVAL}\item [TWAIFD\_BIT\_VAL\_C\_VAL] 配置过滤器 C 的位值。标识符格式与 TX 或 RX buffer 的 IDENTIFIER\_W 相同。如果过滤器 C 不存在，写入此寄存器无效，读取将返回全零。(R/W)
\end{reglist}\end{regdesc}
\end{register}


\begin{register}{H}{TWAIFD\_FILTER\_RAN\_LOW\_REG}{0x0054}\label{regdesc:TWAIFDFILTERRANLOWREG}
\regfield{(reserved)}{3}{29}{000}%
\regfield{TWAIFD\_BIT\_RAN\_LOW\_VAL}{29}{0}{{0x000000}}%
\reglabel{Reset}\regnewline%
\begin{regdesc}\begin{reglist}
\label{fielddesc:TWAIFDBITRANLOWVAL}\item [TWAIFD\_BIT\_RAN\_LOW\_VAL] 配置范围过滤器的阈值下限。标识符格式与 TX 或 RX buffer 的 IDENTIFIER\_W 相同。如果不支持范围过滤器，写入此寄存器无效，读取将返回全零。(R/W)
\end{reglist}\end{regdesc}
\end{register}


\begin{register}{H}{TWAIFD\_FILTER\_RAN\_HIGH\_REG}{0x0058}\label{regdesc:TWAIFDFILTERRANHIGHREG}
\regfield{(reserved)}{3}{29}{000}%
\regfield{TWAIFD\_BIT\_RAN\_HIGH\_VAL}{29}{0}{{0x000000}}%
\reglabel{Reset}\regnewline%
\begin{regdesc}\begin{reglist}
\label{fielddesc:TWAIFDBITRANHIGHVAL}\item [TWAIFD\_BIT\_RAN\_HIGH\_VAL] 配置范围过滤器的阈值上限。标识符格式与 TX 或 RX buffer 的 IDENTIFIER\_W 相同。如果不支持范围过滤器，写入此寄存器无效，读取将返回全零。(R/W)
\end{reglist}\end{regdesc}
\end{register}


\begin{register}{H}{TWAIFD\_FILTER\_CONTROL\_FILTER\_STATUS\_REG}{0x005C}\label{regdesc:TWAIFDFILTERCONTROLFILTERSTATUSREG}
\regfield{(reserved)}{12}{20}{000000000000}%
\regfield{TWAIFD\_SFR}{1}{19}{{1}}%
\regfield{TWAIFD\_SFC}{1}{18}{{1}}%
\regfield{TWAIFD\_SFB}{1}{17}{{1}}%
\regfield{TWAIFD\_SFA}{1}{16}{{1}}%
\regfield{TWAIFD\_FRFE}{1}{15}{{0}}%
\regfield{TWAIFD\_FRFB}{1}{14}{{0}}%
\regfield{TWAIFD\_FRNE}{1}{13}{{0}}%
\regfield{TWAIFD\_FRNB}{1}{12}{{0}}%
\regfield{TWAIFD\_FCFE}{1}{11}{{0}}%
\regfield{TWAIFD\_FCFB}{1}{10}{{0}}%
\regfield{TWAIFD\_FCNE}{1}{9}{{0}}%
\regfield{TWAIFD\_FCNB}{1}{8}{{0}}%
\regfield{TWAIFD\_FBFE}{1}{7}{{0}}%
\regfield{TWAIFD\_FBFB}{1}{6}{{0}}%
\regfield{TWAIFD\_FBNE}{1}{5}{{0}}%
\regfield{TWAIFD\_FBNB}{1}{4}{{0}}%
\regfield{TWAIFD\_FAFE}{1}{3}{{1}}%
\regfield{TWAIFD\_FAFB}{1}{2}{{1}}%
\regfield{TWAIFD\_FANE}{1}{1}{{1}}%
\regfield{TWAIFD\_FANB}{1}{0}{{1}}%
\reglabel{Reset}\regnewline%
\begin{regdesc}\begin{reglist}
\label{fielddesc:TWAIFDFANB}\item [TWAIFD\_FANB] 配置过滤器 A 是否接收 CAN 基本帧。\\
0：未接收\\
1：已接收\\(R/W) 
\label{fielddesc:TWAIFDFANE}\item [TWAIFD\_FANE] 配置过滤器 A 是否接收 CAN 扩展帧。\\
0：未接收\\
1：已接收\\(R/W)
\label{fielddesc:TWAIFDFAFB}\item [TWAIFD\_FAFB] 配置过滤器 A 是否接收 CAN FD 基本帧。\\
0：未接收\\
1：已接收\\(R/W)
\label{fielddesc:TWAIFDFAFE}\item [TWAIFD\_FAFE] 配置过滤器 A 是否接收 CAN FD 扩展帧。\\
0：未接收\\
1：已接收\\(R/W)
\label{fielddesc:TWAIFDFBNB}\item [TWAIFD\_FBNB] 配置过滤器 B 是否接收 CAN 基本帧。\\
0：未接收\\
1：已接收\\(R/W)
\label{fielddesc:TWAIFDFBNE}\item [TWAIFD\_FBNE] 配置过滤器 B 是否接收 CAN 扩展帧。\\
0：未接收\\
1：已接收\\(R/W)
\label{fielddesc:TWAIFDFBFB}\item [TWAIFD\_FBFB] 配置过滤器 B 是否接收 CAN FD 基本帧。\\
0：未接收\\
1：已接收\\(R/W)
\label{fielddesc:TWAIFDFBFE}\item [TWAIFD\_FBFE] 配置过滤器 B 是否接收 CAN FD 扩展帧。\\
0：未接收\\
1：已接收\\(R/W)

\item [见下页...]
\end{reglist}\end{regdesc}
\end{register}

\addtocounter{Regfloat}{-1}
\begin{register}{H}{TWAIFD\_FILTER\_CONTROL\_FILTER\_STATUS\_REG}{0x005C}
\begin{regdesc}\begin{reglist}
\item [...接上页]

\label{fielddesc:TWAIFDFCNB}\item [TWAIFD\_FCNB] 配置过滤器 C 是否接收 CAN 基本帧。\\
0：未接收\\
1：已接收\\(R/W)
\label{fielddesc:TWAIFDFCNE}\item [TWAIFD\_FCNE] 配置过滤器 C 是否接收 CAN 扩展帧。\\
0：未接收\\
1：已接收\\(R/W)
\label{fielddesc:TWAIFDFCFB}\item [TWAIFD\_FCFB] 配置过滤器 C 是否接收 CAN FD 基本帧。\\
0：未接收\\
1：已接收\\(R/W)
\label{fielddesc:TWAIFDFCFE}\item [TWAIFD\_FCFE] 配置过滤器 C 是否接收 CAN FD 扩展帧。\\
0：未接收\\
1：已接收\\(R/W)
\label{fielddesc:TWAIFDFRNB}\item [TWAIFD\_FRNB] 配置范围过滤器是否接收 CAN 基本帧。\\
0：未接收\\
1：已接收\\(R/W)
\label{fielddesc:TWAIFDFRNE}\item [TWAIFD\_FRNE] 配置范围过滤器是否接收 CAN 扩展帧。\\
0：未接收\\
1：已接收\\(R/W)
\label{fielddesc:TWAIFDFRFB}\item [TWAIFD\_FRFB] 配置范围过滤器是否接收 CAN FD 基本帧。\\
0：未接收\\
1：已接收\\(R/W)
\label{fielddesc:TWAIFDFRFE}\item [TWAIFD\_FRFE] 配置范围过滤器是否接收 CAN FD 扩展帧。\\
0：未接收\\
1：已接收\\(R/W)
\label{fielddesc:TWAIFDSFA}\item [TWAIFD\_SFA] 表示过滤器 A 是否可用。\\0：不可用\\1：可用\\(RO)

\item [见下页...]
\end{reglist}\end{regdesc}
\end{register}

\addtocounter{Regfloat}{-1}
\begin{register}{H}{TWAIFD\_FILTER\_CONTROL\_FILTER\_STATUS\_REG}{0x005C}
\begin{regdesc}\begin{reglist}
\item [...接上页]

\label{fielddesc:TWAIFDSFB}\item [TWAIFD\_SFB] 表示过滤器 B 是否可用。\\0：不可用\\1：可用\\(RO)
\label{fielddesc:TWAIFDSFC}\item [TWAIFD\_SFC] 表示过滤器 C 是否可用。\\0：不可用\\1：可用\\(RO)
\label{fielddesc:TWAIFDSFR}\item [TWAIFD\_SFR] 表示范围过滤器是否可用。\\0：不可用\\1：可用\\(RO)
\end{reglist}\end{regdesc}
\end{register}


\begin{register}{H}{TWAIFD\_RX\_DATA}{0x006C}\label{regdesc:TWAIFDRXDATAREG}
\regfield{TWAIFD\_RX\_DATA}{32}{0}{{0}}%
\reglabel{Reset}\regnewline%
\begin{regdesc}\begin{reglist}
\label{fielddesc:TWAIFDRXDATA}\item [TWAIFD\_RX\_DATA] 表示 FIFO 中读取指针位置的 RX buffer 数据。每次读取此寄存器后，只要 RX buffer 中还有下一个数据字，读取指针会自动增加。缓冲区存储的第一个字为 FRAME\_FORMAT\_W，接下来是 IDENTIFIER\_W，依次类推。此寄存器应通过 32 位访问方式读取。(RO)
\end{reglist}\end{regdesc}
\end{register}


\begin{register}{H}{TWAIFD\_TX\_COMMAND\_TXTB\_INFO\_REG}{0x0074}\label{regdesc:TWAIFDTXCOMMANDTXTBINFOREG}
\regfield{(reserved)}{12}{20}{000000000000}%
\regfield{TWAIFD\_TXT\_BUFFER\_COUNT}{4}{16}{{0x4}}%
\regfield{TWAIFD\_TXB8}{1}{15}{{0}}%
\regfield{TWAIFD\_TXB7}{1}{14}{{0}}%
\regfield{TWAIFD\_TXB6}{1}{13}{{0}}%
\regfield{TWAIFD\_TXB5}{1}{12}{{0}}%
\regfield{TWAIFD\_TXB4}{1}{11}{{0}}%
\regfield{TWAIFD\_TXB3}{1}{10}{{0}}%
\regfield{TWAIFD\_TXB2}{1}{9}{{0}}%
\regfield{TWAIFD\_TXB1}{1}{8}{{0}}%
\regfield{(reserved)}{5}{3}{00000}%
\regfield{TWAIFD\_TXCA}{1}{2}{{0}}%
\regfield{TWAIFD\_TXCR}{1}{1}{{0}}%
\regfield{TWAIFD\_TXCE}{1}{0}{{0}}%
\reglabel{Reset}\regnewline%
\begin{regdesc}\begin{reglist}
\label{fielddesc:TWAIFDTXCE}\item [TWAIFD\_TXCE] 配置是否发出 set empty 命令。\\
0：不发出\\
1：发出\\(WO) 
\label{fielddesc:TWAIFDTXCR}\item [TWAIFD\_TXCR] 配置是否发出 set ready 命令。\\
0：不发出\\
1：发出\\(WO)
\label{fielddesc:TWAIFDTXCA}\item [TWAIFD\_TXCA] 配置是否发出 set abort 命令。\\
0：不发出\\
1：发出\\(WO)
\label{fielddesc:TWAIFDTXB1}\item [TWAIFD\_TXB1] 配置是否发出命令到 TX buffer 1。\\
0：不发出\\
1：发出\\(WO)
\label{fielddesc:TWAIFDTXB2}\item [TWAIFD\_TXB2] 配置是否发出命令到 TX buffer 2。\\
0：不发出\\
1：发出\\(WO)
\label{fielddesc:TWAIFDTXB3}\item [TWAIFD\_TXB3] 配置是否发出命令到 TX buffer 3。如果 TX buffer 数量少于 3，此字段保留且无功能。\\
0：不发出\\
1：发出\\(WO)
\label{fielddesc:TWAIFDTXB4}\item [TWAIFD\_TXB4] 配置是否发出命令到 TX buffer 4。如果 TX buffer 数量少于 4，此字段保留且无功能。\\
0：不发出\\
1：发出\\(WO)

\item [见下页...]
\end{reglist}\end{regdesc}
\end{register}

\addtocounter{Regfloat}{-1}
\begin{register}{H}{TWAIFD\_TX\_COMMAND\_TXTB\_INFO\_REG}{0x0074}
\begin{regdesc}\begin{reglist}
\item [...接上页]

\label{fielddesc:TWAIFDTXB5}\item [TWAIFD\_TXB5] 配置是否发出命令到 TX buffer 5。如果 TX buffer 数量少于 5，此字段保留且无功能。\\
0：不发出\\
1：发出\\(WO)
\label{fielddesc:TWAIFDTXB6}\item [TWAIFD\_TXB6] 配置是否发出命令到 TX buffer 6。如果 TX buffer 数量少于 6，此字段保留且无功能。\\
0：不发出\\
1：发出\\(WO)
\label{fielddesc:TWAIFDTXB7}\item [TWAIFD\_TXB7] 配置是否发出命令到 TX buffer 7。如果 TX buffer 数量少于 7，此字段保留且无功能。\\
0：不发出\\
1：发出\\(WO)
\label{fielddesc:TWAIFDTXB8}\item [TWAIFD\_TXB8] 配置是否发出命令到 TX buffer 8。如果 TX buffer 数量少于 8，此字段保留且无功能。\\
0：不发出\\
1：发出\\(WO)
\label{fielddesc:TWAIFDTXTBUFFERCOUNT}\item [TWAIFD\_TXT\_BUFFER\_COUNT] 表示 \modulename 中存在的 TX buffer 数量。最低缓冲区始终为 1。最高缓冲区的索引等于存在的缓冲区数量。(RO)
\end{reglist}\end{regdesc}
\end{register}


\begin{register}{H}{TWAIFD\_TX\_PRIORITY\_REG}{0x0078}\label{regdesc:TWAIFDTXPRIORITYREG}
\regfield{(reserved)}{1}{31}{0}%
\regfield{TWAIFD\_TXT8P}{3}{28}{{0x0}}%
\regfield{(reserved)}{1}{27}{0}%
\regfield{TWAIFD\_TXT7P}{3}{24}{{0x0}}%
\regfield{(reserved)}{1}{23}{0}%
\regfield{TWAIFD\_TXT6P}{3}{20}{{0x0}}%
\regfield{(reserved)}{1}{19}{0}%
\regfield{TWAIFD\_TXT5P}{3}{16}{{0x0}}%
\regfield{(reserved)}{1}{15}{0}%
\regfield{TWAIFD\_TXT4P}{3}{12}{{0x0}}%
\regfield{(reserved)}{1}{11}{0}%
\regfield{TWAIFD\_TXT3P}{3}{8}{{0x0}}%
\regfield{(reserved)}{1}{7}{0}%
\regfield{TWAIFD\_TXT2P}{3}{4}{{0x0}}%
\regfield{(reserved)}{1}{3}{0}%
\regfield{TWAIFD\_TXT1P}{3}{0}{{0x1}}%
\reglabel{Reset}\regnewline%
\begin{regdesc}\begin{reglist}
\label{fielddesc:TWAIFDTXT1P}\item [TWAIFD\_TXT1P] 配置 TX buffer 1 的优先级，数值越高，优先级越高。(R/W) 
\label{fielddesc:TWAIFDTXT2P}\item [TWAIFD\_TXT2P] 配置 TX buffer 2 的优先级，数值越高，优先级越高。(R/W)
\label{fielddesc:TWAIFDTXT3P}\item [TWAIFD\_TXT3P] 配置 TX buffer 3 的优先级，数值越高，优先级越高。如果 TX buffer 数量少于 3，此字段保留且无功能。(R/W)
\label{fielddesc:TWAIFDTXT4P}\item [TWAIFD\_TXT4P] 配置 TX buffer 4 的优先级，数值越高，优先级越高。如果 TX buffer 数量少于 4，此字段保留且无功能。(R/W)
\label{fielddesc:TWAIFDTXT5P}\item [TWAIFD\_TXT5P] 配置 TX buffer 5 的优先级，数值越高，优先级越高。如果 TX buffer 数量少于 5，此字段保留且无功能。(R/W)
\label{fielddesc:TWAIFDTXT6P}\item [TWAIFD\_TXT6P] 配置 TX buffer 6 的优先级，数值越高，优先级越高。如果 TX buffer 数量少于 6，此字段保留且无功能。(R/W)
\label{fielddesc:TWAIFDTXT7P}\item [TWAIFD\_TXT7P] 配置 TX buffer 7 的优先级，数值越高，优先级越高。如果 TX buffer 数量少于 7，此字段保留且无功能。(R/W)
\label{fielddesc:TWAIFDTXT8P}\item [TWAIFD\_TXT8P] 配置 TX buffer 8 的优先级，数值越高，优先级越高。如果 TX buffer 数量少于 8，此字段保留且无功能。(R/W)
\end{reglist}\end{regdesc}
\end{register}


\begin{register}{H}{TWAIFD\_RX\_FR\_CTR\_REG}{0x0084}\label{regdesc:TWAIFDRXFRCTRREG}
\regfield{TWAIFD\_RX\_FR\_CTR\_VAL}{32}{0}{{0}}%
\reglabel{Reset}\regnewline%
\begin{regdesc}\begin{reglist}
\label{fielddesc:TWAIFDRXFRCTRVAL}\item [TWAIFD\_RX\_FR\_CTR\_VAL] 表示接收的帧数。(RO)
\end{reglist}\end{regdesc}
\end{register}


\begin{register}{H}{TWAIFD\_TX\_FR\_CTR\_REG}{0x0088}\label{regdesc:TWAIFDTXFRCTRREG}
\regfield{TWAIFD\_TX\_CTR\_VAL}{32}{0}{{0}}%
\reglabel{Reset}\regnewline%
\begin{regdesc}\begin{reglist}
\label{fielddesc:TWAIFDTXCTRVAL}\item [TWAIFD\_TX\_CTR\_VAL] 表示发送的帧数。(RO)
\end{reglist}\end{regdesc}
\end{register}


\begin{register}{H}{TWAIFD\_DEBUG\_REG}{0x008C}\label{regdesc:TWAIFDDEBUGREG}
\regfield{(reserved)}{13}{19}{0000000000000}%
\regfield{TWAIFD\_PC\_SOF}{1}{18}{{0}}%
\regfield{TWAIFD\_PC\_OVR}{1}{17}{{0}}%
\regfield{TWAIFD\_PC\_SUSP}{1}{16}{{0}}%
\regfield{TWAIFD\_PC\_INT}{1}{15}{{0}}%
\regfield{TWAIFD\_PC\_EOF}{1}{14}{{0}}%
\regfield{TWAIFD\_PC\_ACKD}{1}{13}{{0}}%
\regfield{TWAIFD\_PC\_ACK}{1}{12}{{0}}%
\regfield{TWAIFD\_PC\_CRCD}{1}{11}{{0}}%
\regfield{TWAIFD\_PC\_CRC}{1}{10}{{0}}%
\regfield{TWAIFD\_PC\_STC}{1}{9}{{0}}%
\regfield{TWAIFD\_PC\_DAT}{1}{8}{{0}}%
\regfield{TWAIFD\_PC\_CON}{1}{7}{{0}}%
\regfield{TWAIFD\_PC\_ARB}{1}{6}{{0}}%
\regfield{TWAIFD\_DESTUFF\_COUNT}{3}{3}{{0}}%
\regfield{TWAIFD\_STUFF\_COUNT}{3}{0}{{0}}%
\reglabel{Reset}\regnewline%
\begin{regdesc}\begin{reglist}
\label{fielddesc:TWAIFDSTUFFCOUNT}\item [TWAIFD\_STUFF\_COUNT] 表示当前的填充位计数，对 8 取模，符合 ISO FD 协议定义。每个 CAN 帧开始时填充计数被重置，在到达 ISO FD 帧的填充计数字段之前，每检测到一个填充位加一，之后该值保持不变，直到下一帧开始。在非 ISO FD 或标准 CAN 中，填充位的计数会一直持续到帧结束。请注意，此字段并未采用 ISO FD 标准中规定的格雷码编码。填充计数仅在控制器主动在总线上收发数据时更新；在接收过程中，此值保持不变。(RO)
\label{fielddesc:TWAIFDDESTUFFCOUNT}\item [TWAIFD\_DESTUFF\_COUNT] 表示当前的去填充位计数，对 8 取模，符合 ISO FD 协议定义。每个帧开始时去填充计数被重置，在到达 ISO FD 帧的填充计数字段之前，每检测到一个去填充位加一，之后该值保持不变，直到下一帧开始。在非 ISO FD 或标准 CAN 中，去填充位的计数会一直持续到帧结束。请注意，此字段并未采用 ISO FD 标准中规定的格雷码编码。去填充计数在发送和接收过程中都会被更新。(RO)
\label{fielddesc:TWAIFDPCARB}\item [TWAIFD\_PC\_ARB] 表示协议控制状态机是否处于仲裁字段。\\0：不处于\\1：处于\\(RO) 
\label{fielddesc:TWAIFDPCCON}\item [TWAIFD\_PC\_CON] 表示协议控制状态机处于控制字段。\\0：不处于\\1：处于\\(RO)
\label{fielddesc:TWAIFDPCDAT}\item [TWAIFD\_PC\_DAT] 表示协议控制状态机处于数据字段。\\0：不处于\\1：处于\\(RO)
\label{fielddesc:TWAIFDPCSTC}\item [TWAIFD\_PC\_STC] 表示协议控制状态机处于填充计数字段。\\0：不处于\\1：处于\\(RO)
\label{fielddesc:TWAIFDPCCRC}\item [TWAIFD\_PC\_CRC] 表示协议控制状态机处于 CRC 字段。\\0：不处于\\1：处于\\(RO)

\item [见下页...]
\end{reglist}\end{regdesc}
\end{register}

\addtocounter{Regfloat}{-1}
\begin{register}{H}{TWAIFD\_DEBUG\_REG}{0x008C}
\begin{regdesc}\begin{reglist}
\item [...接上页]

\label{fielddesc:TWAIFDPCCRCD}\item [TWAIFD\_PC\_CRCD] 表示协议控制状态机处于 CRC 分隔符字段。\\0：不处于\\1：处于\\(RO)
\label{fielddesc:TWAIFDPCACK}\item [TWAIFD\_PC\_ACK] 表示协议控制状态机处于 ACK 字段。\\0：不处于\\1：处于\\(RO)
\label{fielddesc:TWAIFDPCACKD}\item [TWAIFD\_PC\_ACKD] 表示协议控制状态机处于 ACK 分隔符字段。\\0：不处于\\1：处于\\(RO)
\label{fielddesc:TWAIFDPCEOF}\item [TWAIFD\_PC\_EOF] 表示协议控制状态机处于文件结束字段。\\0：不处于\\1：处于\\(RO)
\label{fielddesc:TWAIFDPCINT}\item [TWAIFD\_PC\_INT] 表示协议控制状态机处于间歇字段。\\0：不处于\\1：处于\\(RO)
\label{fielddesc:TWAIFDPCSUSP}\item [TWAIFD\_PC\_SUSP] 表示协议控制状态机处于暂停发送字段。\\0：不处于\\1：处于\\(RO)
\label{fielddesc:TWAIFDPCOVR}\item [TWAIFD\_PC\_OVR] 表示协议控制状态机处于过载字段。\\0：不处于\\1：处于\\(RO)
\label{fielddesc:TWAIFDPCSOF}\item [TWAIFD\_PC\_SOF] 表示协议控制状态机处于帧开始字段。\\0：不处于\\1：处于\\(RO)
\end{reglist}\end{regdesc}
\end{register}


\begin{register}{H}{TWAIFD\_DATE\_VER\_REG}{0x0FFC}\label{regdesc:TWAIFDDATEVERREG}
\regfield{TWAIFD\_DATE\_VER}{32}{0}{{0x2312150}}%
\reglabel{Reset}\regnewline%
\begin{regdesc}\begin{reglist}
\label{fielddesc:TWAIFDDATEVER}\item [TWAIFD\_DATE\_VER] 版本控制寄存器。(R/W)
\end{reglist}\end{regdesc}
\end{register}
